\documentclass[10pt,a4paper]{amsart}
\usepackage[margin=19mm]{geometry}
\usepackage{amsmath,amssymb,mathtools}
\usepackage[scr=rsfs,cal=euler]{mathalfa}
\usepackage{xfrac}
\usepackage[unicode,hidelinks,bookmarks=false]{hyperref}
\newcommand{\ie}{i.e.\ }
\newcommand{\cf}{cf.\ }
\newcommand{\resp}{respectively\ }
\newcommand{\Subsection}{\subsection}
\providecommand{\nobreakdash}{\nobreak\hbox{-}}
\setlength{\emergencystretch}{2em}
\multlinegap0pt
\usepackage{algorithmic}
\usepackage{float}
\usepackage{enumitem}
\setlist[enumerate]{leftmargin=.5in}
\setlist[itemize]{leftmargin=.5in}
\usepackage{bm}
\usepackage{amsopn}
\usepackage{aliascnt}
\usepackage[capitalize,nameinlink]{cleveref}
\makeatletter
\Crefname{ALC@unique}{Line}{Lines}
\makeatother
\crefname{hypothesis}{Hypothesis}{Hypotheses}
\crefname{algorithm}{Algorithm}{Algorithms}
\crefname{section}{section}{sections}
\crefname{subsection}{subsection}{subsections}
\Crefname{section}{Section}{Sections}
\Crefname{subsection}{Subsection}{Subsections}
\Crefname{figure}{Figure}{Figures}
\crefformat{equation}{\textup{#2(#1)#3}}
\crefrangeformat{equation}{\textup{#3(#1)#4--#5(#2)#6}}
\crefmultiformat{equation}{\textup{#2(#1)#3}}{ and \textup{#2(#1)#3}}{, \textup{#2(#1)#3}}{, and \textup{#2(#1)#3}}
\crefrangemultiformat{equation}{\textup{#3(#1)#4--#5(#2)#6}}{ and \textup{#3(#1)#4--#5(#2)#6}}{, \textup{#3(#1)#4--#5(#2)#6}}{, and \textup{#3(#1)#4--#5(#2)#6}}
\Crefformat{equation}{#2Equation~\textup{(#1)}#3}
\Crefrangeformat{equation}{Equations~\textup{#3(#1)#4--#5(#2)#6}}
\Crefmultiformat{equation}{Equations~\textup{#2(#1)#3}}{ and \textup{#2(#1)#3}}{, \textup{#2(#1)#3}}{, and \textup{#2(#1)#3}}
\Crefrangemultiformat{equation}{Equations~\textup{#3(#1)#4--#5(#2)#6}}{ and \textup{#3(#1)#4--#5(#2)#6}}{, \textup{#3(#1)#4--#5(#2)#6}}{, and \textup{#3(#1)#4--#5(#2)#6}}
\crefdefaultlabelformat{#2\textup{#1}#3}
\numberwithin{equation}{section}
\floatstyle{ruled}
\newfloat{algorithm}{htbp}{loa}
\floatname{algorithm}{Algorithm}
\numberwithin{algorithm}{section}
\newcommand{\creflastconjunction}{, and~}
\theoremstyle{plain}
\newtheorem{theorem}{Theorem}[section]
\newaliascnt{lemma}{theorem}
\newtheorem{lemma}[lemma]{Lemma}
\aliascntresetthe{lemma}
\newaliascnt{corollary}{theorem}
\newtheorem{corollary}[corollary]{Corollary}
\aliascntresetthe{corollary}
\newaliascnt{proposition}{theorem}
\newtheorem{proposition}[proposition]{Proposition}
\aliascntresetthe{proposition}
\newaliascnt{definition}{theorem}
\newtheorem{definition}[definition]{Definition}
\aliascntresetthe{definition}
\newaliascnt{claim}{theorem}
\newtheorem{claim}[claim]{Claim}
\aliascntresetthe{claim}
\newaliascnt{assumption}{theorem}
\newtheorem{assumption}[assumption]{Assumption}
\aliascntresetthe{assumption}
\theoremstyle{remark}
\newaliascnt{remark}{theorem}
\newtheorem{remark}[remark]{Remark}
\aliascntresetthe{remark}
\newaliascnt{hypothesis}{theorem}
\newtheorem{hypothesis}[hypothesis]{Hypothesis}
\aliascntresetthe{hypothesis}
\renewcommand{\algorithmicrequire}{\textbf{Input:}}
\renewcommand{\algorithmicensure}{\textbf{Output:}}
\DeclareMathOperator*{\argmin}{argmin}
\DeclareMathOperator{\diverge}{div}
\DeclareMathOperator{\supp}{supp}
\DeclareMathOperator{\var}{var}
\DeclareMathOperator{\cov}{cov}
\DeclareMathOperator{\arctanh}{arctanh}
\DeclareMathOperator{\cond}{cond}
\DeclareMathOperator{\diag}{diag}
\newcommand{\C}{\mathbb{C}}
\newcommand{\op}{\operatorname}
\newcommand{\Q}{\mathbb{Q}}
\newcommand{\R}{\mathbb{R}}
\newcommand{\N}{\mathbb{N}}
\newcommand{\bE}{\mathbb{E}}
\newcommand{\dd}{\mathrm{d}}
\newcommand*\Bell{\ensuremath{\boldsymbol\ell}}
\renewcommand{\Tilde}[1]{\widetilde{#1}}
\title[Global convergence of gradient-enhanced least squares Monte Carlo]{Gradient-enhanced sparse Hermite polynomial expansions for pricing and hedging high-dimensional American options: Theorem 5.6}
\author{Jiefei Yang, Guanglian Li}
\date{Source: arXiv:2405.02570v2 (CC BY 4.0)}
\begin{document}
\maketitle
{\sloppy\noindent\emph{Source and licence.} Gradient-enhanced sparse Hermite polynomial expansions for pricing and hedging high-dimensional American options. Version arXiv:2405.02570v2 (CC BY 4.0). This material is distributed under the Creative Commons Attribution 4.0 International licence, http://creativecommons.org/licenses/by/4.0/, as recorded in the arXiv OAI metadata for this exact version and shown on the abstract page https://arxiv.org/abs/2405.02570v2. Full rights notice: licenses/arxiv-2405.02570-CC-BY-4.0.txt.\par}\smallskip
\noindent\textit{Editorial scope.} Counted claim: original Theorem 5.6 (source label 'thm: global error', source0.tex lines 691-697), the main theorem of the article ("we shall prove the global convergence of G-LSM ... c.f. Theorem 5.6"), delivered with its complete original proof at source lines 698-708 as the single primary proof body. Required context retained uncounted: Sections 1 to 5 of the article as printed (source lines 118-690), that is the introduction and notation, the Bermudan pricing and hedging framework with the backward induction (2.1) and the continuation value function, the sparse Hermite polynomial expansion with the hyperbolic cross index set (3.2) and the derivative formula (3.6), the gradient-enhanced least squares and G-LSM algorithms 4.1 and 4.2, the stability Theorem 4.4 with its proof, Remarks 4.2, 4.3 and 4.5, Assumption 5.1, Lemma 5.2, Lemma 5.3, Theorem 5.4 and Corollary 5.5 with their complete proofs, and the paragraph that follows the counted proof (source lines 710-711). Printed numbering is reproduced with the author's own declared environments (theorem, lemma, corollary, proposition, definition, remark, hypothesis, claim and assumption carried by alias counters of one theorem counter numbered within the section, algorithms numbered within the section, equations numbered within the section, and the figure numbered continuously as the siamonline220329 option onefignum requires, with the class's own cleveref formats for equations, sections and figures), so that the delivered PDF prints Theorem 4.1, Remark 4.2, Remark 4.3, Theorem 4.4, Remark 4.5, Assumption 5.1, Lemma 5.2, Lemma 5.3, Theorem 5.4, Corollary 5.5, Theorem 5.6, Algorithm 4.1, Algorithm 4.2 and Figure 1 exactly as the published PDF does. The publisher's siamonline220329 class is not present in the runtime; the unit is delivered inside the AMS amsart wrapper with the author's own macros, the algorithmic package for the two algorithm bodies, a ruled float shim that restores the class's Algorithm environment, and an editorial alias that renders the author's \textbackslash{}Tilde accent used in the notation paragraph. Omissions declared under a bounded\_source\_comparison receipt: the three figure-artwork lines \textbackslash{}includegraphics[width = .3\textbackslash{}textwidth]{hc\_nnz\_order4.eps}, {hc\_nnz\_order6.eps} and {hc\_nnz\_order10.eps} inside Figure 1 (the figure environment, its tabular layout, its subcaptions, its caption and its label fig: hc\_nnz are retained and the figure is still cross-referenced, but no artwork is redistributed), and three navigational sentences that point at the sections outside the unit (the closing sentences of the structure paragraph at source line 136, "Their numerical performance will be compared in \textbackslash{}cref{sec:examples}." at source line 355, "We will numerically verify the complexity analysis in \textbackslash{}cref{subsec:ex3}." at source line 398 and "See \textbackslash{}cref{subsec:heston} for one numerical example related to the Heston model." at source line 519). Slips kept literal: the stray closing parenthesis in the estimate E[(\textbackslash{}check{Z}\_k - \textbackslash{}tilde{Z}\_{t\_k})\textasciicircum{}2)] of the proof of Lemma 5.3 (source line 597) and the inconsistent notation c\_k\textasciicircum{}{CLS} against c\_k\textasciicircum{}{\textbackslash{}op{CLS}} in the proof of Corollary 5.5 (source line 673) are reproduced without change. Reference handling: the delivered unit keeps every \textbackslash{}cite command of the retained spans and transcribes exactly the cited entries, adcock2022sparse, bayer2023pricing, becker2019deep, becker2020pricing, bouchard2012monte, chen2021deep, E2017deepbsde, el1997reflected, gao2023convergence, gobet2007error, gobet2005regression, guyon2013nonlinear, Hestenes:1952, hure2020deep, lapeyre2021neural, longstaff2001valuing, ludkovski2018kriging, luo2018error, oksendal2013stochastic, seydel2006tools, sirignano2018dgm, tsitsiklis2001regression, wang2018deep, wang2009pricing, yang2023optionpricing, zhang2017backward, from the article's own bibliography with the published (siamplain) numbers as their labels, so that every citation prints the same number as the published PDF. No mathematics is written, repaired or re-derived; all mathematics is native text.
\section{Introduction}
The early exercise feature of American or Bermudan options gives holders the right to buy (call) or sell (put) underlying assets before the expiration date, and their accurate numerical calculation is of great practical importance. Meanwhile, the efficient estimation of Greeks (i.e., derivatives of the price function, e.g., delta and gamma) is vital for hedging and risk management, since the theory of option pricing builds on the assumption of the absence of arbitrage. For example, when the asset price is on the rise, the gain in the long position of a call writer's asset may offset the potential loss of the call option.

Nonetheless, the early exercise feature of American options poses significant challenges for computing price and Greeks, especially in high dimensions. One of the most popular methods for high-dimensional American option pricing is the least squares Monte Carlo (LSM) method \cite{longstaff2001valuing, tsitsiklis2001regression}. Computing Greeks in high dimensions is more involved, and further developments with LSM have been proposed \cite{wang2009pricing, bouchard2012monte}.
In this work, building on LSM, we shall develop a simple, fast, and accurate algorithm, termed as gradient-enhanced least squares Monte Carlo (G-LSM) method, cf. \cref{alg: sparse hermite polynomial expansions}, for computing price and Greeks simultaneously at all time steps for dimensions up to 100. The key methodological innovations includes using sparse Hermite polynomial space with a hyperbolic cross index set as the ansatz space for approximating the continuation value functions (CVFs), and incorporating the gradient information for computing the expansion coefficients.

Next we elaborate the two methodological innovations. First, the main obstacle of using gradient information lies in the computational expense: a $d$-variate price has $d$ partial derivatives, which grows quickly with $d$, especially when the derivative evaluation is costly. With the proposed sparse Hermite polynomial ansatz space, the derivatives of polynomial bases can be obtained at almost no extra cost, cf. \cref{eq:derivative_herm} below. This allows greatly reducing the computational cost. Second, although using a polynomial ansatz space for the CVFs as LSM, G-LSM finds  the expansion coefficients via solving a linear least squares problem enhanced by the gradient (and hence the name G-LSM), by minimizing the mean squared error between the approximate and exact value functions at $t_{k+1}$, cf. \cref{eq: empirical risk} below. This differs markedly from LSM, which approximates the conditional expectations by projection and minimizes the mean squared error of approximating CVF at $t_k$. Numerical experiments show that this choice can achieve better accuracy in price, Greeks, and optimal exercise strategies than LSM. Although other classes of basis functions, e.g., Chebyshev, Legendre, and Laguerre polynomials, might be applicable, we propose to adopt sparse Hermite polynomials as the ansatz space for CVF because (1) CVF can be reformulated as a function of independent Brownian motions with its best approximation error ensured by  \cref{lem:best-error-HC} and (2) the gradient of Hermite polynomials can be evaluated very efficiently, which is the key point to enable the algorithm to work in high dimensions.

In G-LSM, the approximation of the terminal condition at $t_{k+1}$ is obtained by discretizing the linear backward stochastic differential equation (BSDE) for the CVF, c.f. \cref{thm: linear bsde}, which was recently innovated in a deep neural network-based method for American option pricing \cite{chen2021deep}. The idea of matching the terminal condition has been widely applied in solving high-dimensional BSDEs with deep neural networks (DNNs) \cite{E2017deepbsde, hure2020deep}. In practice, it involves computing the gradient of the CVF, and in turn that of the basis functions in the ansatz space (in addition to function evaluation). When $d\gg1$, for a complicated ansatz space, evaluating the derivatives at all time steps can be prohibitive. We shall show that the extra cost is nearly negligible for the sparse Hermite polynomial ansatz, and that the overall complexity of G-LSM with $N$ time steps, $M$ sample paths, and $N_b$ basis functions is $\mathcal{O}(NMN_b)$, nearly identical with that for LSM. Numerical results show that the accuracy of G-LSM is competitive with DNN-based methods for dimensions up to $d = 100$.

In theory, the CVF can be formulated as a smooth, high-dimensional function in ${L}^2_{\omega}(\R^d)$ with a Gaussian weight function $\omega(\mathbf{y})$ \cite{yang2023optionpricing}.
This regularity enables the use of normalized and generalized Hermite polynomials, which form an orthonormal basis of $L^2_{\omega}(\R^d)$. Furthermore, drawing on the geometric convergence rate of the hyperbolic cross approximation with Hermite polynomials \cite{luo2018error}, we shall prove the global convergence of G-LSM using BSDE technique, stochastic and Malliavin calculus, and establish an error bound in terms of time step size, statistical error of the Monte Carlo approximation, and the best approximation error in weighted Sobolev space $H^1_{\omega}(\R^d)$, c.f. \cref{thm: global error}. In sum, the algorithmic development of G-LSM, its error analysis and extensive numerical evaluation represent the main contributions of the present work.
% \textcolor{red}{In contrast, the convergence of the DNN-based methods \cite{gao2023convergence,hure2020deep} relies on the assumption that the solution obtained from the neural network approximation is not stuck in a local minimum. This assumption is theoretically unsound due to the non-convex nature of the corresponding loss function.} \noteYang{I suggest that we should not criticize the DNN-based methods in the paper.} 

Now we situate the present study in existing works. Currently, there are two popular classes of methods to price American options in high dimensions: (i) least-squares Monte Carlo-based (LSM) methods and (ii) DNN-based methods. The LSM method has shown tremendous success for pricing American or Bermudan options with more than one stochastic factors. The original LSM \cite{longstaff2001valuing} uses polynomials to approximate the CVF, and other choices have also been explored, e.g., Gaussian process \cite{ludkovski2018kriging} and DNNs \cite{lapeyre2021neural, becker2020pricing}. Recently, LSM with the hierarchical tensor train technique has been studied in \cite{bayer2023pricing}, which demonstrates the success of polynomial approximation for CVFs in very high dimensions. The proposed G-LSM is a variant of LSM that incorporates gradient information that comes nearly for free. Due to the excellent capability for high-dimensional approximation of DNNs, several methods based on DNNs have been proposed for pricing American or Bermudan options, based on optimal stopping problem (parameterizing the stopping time by DNNs and then maximizing the expected reward \cite{becker2019deep}), free boundary PDEs (parameterizing PDE solutions with DNNs \cite{sirignano2018dgm}), or BSDEs (parameterizing the solution pair of the associated reflected BSDE \cite{el1997reflected} by DNNs \cite{hure2020deep}). Within the framework of BSDEs, Chen and Wan \cite{chen2021deep} suggest approximating the difference of the CVF between adjacent time steps by averaging several trained neural networks, which has a quadratic complexity in the number of time steps. Wang et al. \cite{wang2018deep} extend the deep BSDE method \cite{E2017deepbsde} from European option pricing to Bermudan one, with the loss function being the variance of the initial value, and Gao et al. \cite{gao2023convergence} analyze its convergence. In comparison with DNN-based methods, G-LSM enjoys high efficiency and robustness, which involves only least-squares problems and is easy to implement.

The proposed G-LSM scheme is closely linked to the regression-based Monte Carlo method for solving BSDEs \cite{gobet2005regression}. It involves solving a sequence of least squares problems backward in time, where a linear combination of general basis functions is used to approximate each component of the solution $(Y_{t_k}, Z_{1, t_k}, \dots, Z_{d, t_k})$ to the BSDE.  However, this scheme may become infeasible for high dimensions \cite{hure2020deep}. In sharp contrast, in the $d$-dimensional case, the G-LSM method has a computational complexity $d$ times less that of the approach in \cite{gobet2005regression}, while maintains the stability of least squares regression. See Remark \ref{rmk:complexity} for further discussions on the complexity and stability. The efficiency of the scheme \cite{gobet2005regression} for high-dimensional problems was enhanced by implementing a neural network in the scheme \cite{hure2020deep}, which has demonstrated impressive empirical performance.

The structure of this article is organized as follows. In \cref{sec:pricing-hedging}, we describe the mathematical framework of pricing and hedging high-dimensional American or Bermudan options, and in \cref{sec:sparse-hermite}, we recall several useful properties of generalized Hermite polynomials and approximation with sparse hyperbolic cross index set. Then in \cref{sec:algorithm-cost}, we derive the main algorithm, i.e., gradient-enhanced least squares Monte Carlo (G-LSM) method, and establish its local and global error estimates in \cref{sec:convergence}.

Throughout, bold and plain letters represent multi-variable and scalars, respectively, and the capital and bold letters $\mathbf{S}$, $\Tilde{\mathbf{X}}$ and $\mathbf{W}$ denote random vectors. The notation $\mathbf{a} \cdot \mathbf{b}$ denotes the dot product of two vectors $\mathbf{a}$ and $\mathbf{b}$, $\operatorname{Tr}(A)$ the trace of a square matrix $A$, and $^\top$ vector transpose. The notation
$\nabla_\mathbf{x}f(t, \mathbf{x})$ and  $\operatorname{Hess}_{\mathbf{x}}f(t,\mathbf{x})$ denote respectively the gradient and Hessian of $f$ with respect to  $\mathbf{x}$. For a multi-index $\bm{\alpha} = (\alpha_1,\ldots,\alpha_d)^\top \in \mathbb{N}_0^d$, we denote $|\bm{\alpha}|: = \alpha_1 + \dots + \alpha_d$. For a positive-valued integrable function $\omega: \mathbb{R}^d\to \mathbb{R}^+$, the weighted space $L^2_{\omega}(\mathbb{R}^d)$ is defined by
\begin{align*}
L^2_{\omega}(\mathbb{R}^d):=\{f:\mathbb{R}^d\to \mathbb{R}: \|f\|_{L^2_{\omega}(\mathbb{R}^d)}<\infty\},\quad \text{ with }  \|f\|_{L^2_{\omega}(\mathbb{R}^d)}^2:=\int_{\mathbb{R}^d}f(\mathbf{x})^2\omega(\mathbf{x})\mathrm{d}\mathbf{x}.
\end{align*}
The weighted Sobolev space $H_{\omega}^m(\mathbb{R}^d)$, $m\in \mathbb{N}$, is defined by
$$
H_{\omega}^m(\mathbb{R}^d) := \{f: \R^d \to \R: \|f\|_{H_\omega^m(\R^d)} < \infty\}, \quad \text{with } \|f\|_{H_\omega^m(\R^d)}^2 = \sum_{0\le |\bm{\alpha}|\le m} \left\|\frac{\partial^{\bm{\alpha}} f}{\partial \mathbf{x}^{\bm{\alpha}}}\right\|_{L^2_\omega(\R^d)}^2.
$$
\section{Bermudan option pricing and hedging} \label{sec:pricing-hedging}
Now we describe the valuation framework for American or Bermudan option pricing and hedging.

\subsection{Option pricing and Greeks}
The fair price of American option $v^A(t)$ at time $t\in [0,T]$ is expressed as the solution to the optimal stopping problem in a risk neutral probability space $(\Omega, \mathcal{F}, (\mathcal{F}_t)_{0\le t\le T},\mathbb{Q})$,
\begin{equation*}
    v^A(t) = \sup_{\tau_t \in [t,T]} \mathbb{E}[e^{-r(\tau_t - t)} g(\mathbf{S}_{\tau_t} ) | \mathcal{F}_t],
\end{equation*}
where $\tau_t$ is an $\mathcal{F}_t$-stopping time, $T>0$ is the expiration date, $(\mathbf{S}_t)_{0\le t\le T}$ is a collection of $d$-dimensional price processes, and $g(\mathbf{S}_t) \in L^2(\Omega, \mathcal{F}_t, \mathbb{Q})$ is the payoff depending on the type of the option.

Numerically, the price of Bermudan option is used to approximate the American one. The Bermudan option can be exercised at finite discrete times $0 = t_0 < t_1 < \dots < t_N = T$ with $\Delta t := t_{k+1} - t_k$ for all $k = 0,1,\dots, N-1$. Using dynamic programming principle or Snell envelope theory \cite[Section 1.8.4]{seydel2006tools}, the Bermudan price function ${v}_{t_k}$ at time $t_k$ is given by the following backward induction:
\begin{equation} \label{eq: backward induction}
    \begin{aligned}
        v_{t_N}(\mathbf{s}) &= g_{t_N}(\mathbf{s}), \\
        v_{t_k}(\mathbf{s}) &= \begin{cases}
            g_{t_k}(\mathbf{s}), \quad \text{ if } g_{t_k}(\mathbf{s}) \ge \tilde{c}_{t_k}(\mathbf{s}), \\
            \tilde{c}_{t_k}(\mathbf{s}), \quad \text{ if } g_{t_k}(\mathbf{s}) < \tilde{c}_{t_k}(\mathbf{s}),
        \end{cases} \quad \text{ for } k = N-1:-1:0,
    \end{aligned}
\end{equation}
where $g_{t_k}$ is the discounted payoff (exercise value) function and $\tilde{c}_{t_k}$ is the CVF, defined by
\begin{equation} \label{eq: continuation value function S}
    \tilde{c}_{t_k}(\mathbf{s}) = \mathbb{E}[v_{t_{k+1}}(\mathbf{S}_{t_{k+1}}) | \mathbf{S}_{t_k} = \mathbf{s}].
\end{equation}
The options delta and gamma are defined to be the first and second order derivatives of the price function $v_{t_k}$ with respect to the price of underlying assets:
\begin{equation} \label{eq: defi greeks}
    \begin{aligned}
        \Delta_{t_k} &:= \nabla v_{t_k}(\mathbf{s}) = \left(\frac{\partial v_{t_k}(\mathbf{s})}{\partial s_1}, \dots, \frac{\partial v_{t_k}(\mathbf{s})}{\partial s_d} \right)^\top\quad \mbox{and}\quad
        \Gamma^{ij}_{t_k} &:= \frac{\partial^2 v_{t_k}(\mathbf{s})}{\partial s_j \partial s_i}.
    \end{aligned}
\end{equation}
We consider all exercise and continuation values of Bermudan option discounted to the time $t = 0$. Since the discrete-time approximation for American option pricing is always needed in practical computation, we will make no semantic difference between American and Bermudan options, following \cite[p. 133]{guyon2013nonlinear}.

\subsection{Multi-asset model and transformation} \label{sec:multi model}
One of the most classical models for high-dimensional American option pricing is the multi-asset Black-Scholes model. Under the risk-neutral probability $\mathbb{Q}$, the prices of $d$ underlying assets, $\mathbf{S}_t = \left(S^1_t, \dots, S^d_t \right)^\top$, follow the correlated geometric Brownian motions
\begin{equation} \label{eq: multi asset BS model}
    \mathrm{d}S^i_t = (r - \delta_i) S^i_t\,\mathrm{d}t + \sigma_i S^i_t\,\mathrm{d}\tilde{W}^i_t \quad \text{ with } S^i_0 = s^i_0, \quad i = 1,2,\dots,d,
\end{equation}
where $\tilde{W}^i_t$ are correlated Brownian motions with correlation $\mathbb{E}[\mathrm{d}\tilde{W}^i_t \mathrm{d}\tilde{W}^j_t] = \rho_{ij} \,\mathrm{d}t$, and $r$, $\delta_i$ and $\sigma_i$ are the riskless interest rate, dividend yields,
and volatility parameters, respectively. We
denote the correlation matrix by $P = (\rho_{ij})_{d\times d}$, the volatility matrix by $\Sigma$ (which is a diagonal matrix with
volatility $\sigma_i$ on the diagonal), and write the dividend yields as a vector $\bm{\delta} = [\delta_1, \dots, \delta_d]^\top$. Using the spectral decomposition $\Sigma P \Sigma^\top=Q\Lambda Q^\top$, the rotated log-price $\Tilde{\mathbf{X}}_t := Q^\top \ln(\mathbf{S}_t./\mathbf{s}_0)$ satisfies an independent Gaussian distribution
\begin{align*}
\Tilde{\mathbf{X}}_t\sim
\mathcal{N}\left(Q^\top \left(r - \bm{\delta} - \frac{1}{2}\Sigma^2\mathbf{1}\right)t, \Lambda t\right).
\end{align*}
Let $\bm{\mu} = Q^\top \left(r - \bm{\delta} - \frac{1}{2}\Sigma^2\mathbf{1}\right)$ and $\lambda_i$ be the $i$-th diagonal element of $\Lambda$. This gives a transformation between underlying asset prices $\mathbf{S}_t$ and independent Brownian motions $\mathbf{W}_t = [W^1_t, \dots, W^d_t]^\top$, i.e.,
\begin{equation} \label{eq: trans S and W}
    \mathbf{S}_t = \mathbf{s}_0 \odot \exp\left(Q \left(\bm{\mu}t + \sqrt{\Lambda} \mathbf{W}_t \right) \right),
\end{equation}
where $\odot$ denotes componentwise product.

From \cref{eq: trans S and W} and \cref{eq: continuation value function S}, the CVF $c_{t_k}$ with respect to the independent Brownian motions is
\begin{equation} \label{eq: continuation value W}
    c_{t_k}(\mathbf{w}) = \mathbb{E}[ u_{t_{k+1}}(\mathbf{W}_{t_{k+1}}) | \mathbf{W}_{t_k} = \mathbf{w}],
\end{equation}
where $u_{t_{k+1}}$ represents the value function at time $t_{k+1}$ with respect to the independent Brownian motions. Our aim is to develop a gradient-enhanced least squares Monte Carlo method that can efficiently approximate $c_{t_k}$, and thus provide accurate prices and their derivatives.

\section{Sparse Hermite polynomial expansion and gradient}\label{sec:sparse-hermite}
Sparse polynomial chaos expansion can serve as a surrogate model of unknown stochastic variables with finite second-order moments. The motivations of using sparse Hermite polynomial expansion for pricing and hedging American options are twofold:
\begin{enumerate}
    \item Let $\omega_t$, $t>0$, be the Gaussian density function defined by $\omega_t(\mathbf{y}) := \prod_{j=1}^d \rho_{t}(y_j)$ with $\rho_t(y):= \frac{1}{\sqrt{2\pi t}} \exp(-\frac{y^2}{2t})$. The CVF $c_{t_k}$ satisfies $c_{t_k} \in {L}^2_{\omega_{t_k}}(\R^d)$, since the payoff $g(\mathbf{S}_{t_k})\in L^2(\Omega, \mathcal{F}_{t_k}, \mathbb{Q})$ has a finite second-order moment. Thus, as an orthonormal basis for $L^2_{\omega_{t_k}}(\R^d)$, the set of normalized and generalized Hermite polynomials is a natural choice.
    \item The CVF $c_{t_k} $ has the regularity ${H}^m_{\omega_{t_k}}(\R^d)$ for any positive integer $m$ by \cite[Lemmas 4.2 and 4.3]{yang2023optionpricing}. The smoothness of the function implies efficient polynomial approximation.
\end{enumerate}

Next, we recall one-dimensional normalized and generalized Hermite polynomials $H_n^{(t)}(y)$ and the tensorized $d$-dimensional polynomials $H_{\bm{\alpha}}^{(t)}(\mathbf{y})$, $\bm{\alpha} \in \mathbb{N}_0^d$. For $d = 1$, using the standard probabilist’s Hermite polynomials $H_n(x) = (-1)^n e^{\frac{x^2}{2}} \frac{\dd^n}{\dd x^n}e^{-\frac{x^2}{2}}$, we define the  $n$th-order normalized and generalized Hermite polynomial $H_n^{(t)}$ by
\begin{equation*}
    H_n^{(t)}(y) := \frac{H_n(\frac{y}{\sqrt{t}})}{\sqrt{n!}}, \quad \text{ for } n\in\mathbb{N}_0, t>0, y\in \R.
\end{equation*}
Then $\{H_n^{(t)}\}_{n\in \mathbb{N}_0}$ forms a complete orthonormal basis for ${L}^2_{\rho_{t}}(\R)$:
\begin{equation*}
    \mathbb{E}\left[ H_n^{(t)}(W_t) H_m^{(t)}(W_t) \right] = \int_\R H_n^{(t)}(y) H_m^{(t)}(y) \rho_t(y) \,\mathrm{d}y = \delta_{nm},
\end{equation*}
with $\delta_{nm}$ being the Kronecker delta.
For $d>1$, the tensorized Hermite polynomial with the multi-index $\bm{\alpha}=(\alpha_1,\ldots,\alpha_d)^\top \in \mathbb{N}_0^d$ defined by
\begin{equation*}
    H_{\bm{\alpha}}^{(t)}(\mathbf{y}) := \prod_{j=1}^d H_{\alpha_j}^{(t)}(y_j)
\end{equation*}
forms a complete orthonormal basis for ${L}^2_{\omega_{t}}(\R^d)$, satisfying
\begin{equation*}
    \mathbb{E}\left[ H_{\bm{\alpha}}^{(t)}(\mathbf{W}_t) H_{\bm{\gamma}}^{(t)}(\mathbf{W}_t) \right] = \int_{\R^d} H_{\bm{\alpha}}^{(t)}(\mathbf{y}) H_{\bm{\gamma}}^{(t)}(\mathbf{y})\omega_{t}(\bm{y}) \,\mathrm{d}\mathbf{y} = \delta_{\bm{\alpha}, \bm{\gamma}}, \quad \bm{\alpha}, \bm{\gamma} \in \mathbb{N}_0^d.
\end{equation*}
Here, $\delta_{\bm{\alpha}, \bm{\gamma}} = \prod_{j=1}^d \delta_{\alpha_j,\gamma_j}$ and
$\omega_{t}(\bm{y})=\prod_{j=1}^d \rho_t(y_j) $.

For any fixed multi-index set $I \subset \mathbb{N}_0^d$, the Hermite polynomial ansatz space $P_{I,k}$ is defined by
\begin{align}\label{eq:pkI}
P_{I,k}:=\text{span}\left\{H_{\bm{\alpha}}^{(t_k)}: \bm{\alpha} \in I\right\}.
\end{align}
Then we aim to approximate the CVF $ c_{t_k}$ in $P_{I,k}$ by
\begin{equation*}
    c_{t_k}(\mathbf{W}_{t_k}) \approx \sum_{\bm{\alpha}\in I} \beta_{\bm{\alpha}} H_{\bm{\alpha}}^{(t_k)}(\mathbf{W}_{t_k}).
\end{equation*}
It is well-known that computing polynomial approximations in high dimensions suffers from the notorious curse of dimensionality with tensor product-type multi-index sets. Fortunately, the smoothness of $c_{t_k}$ implies a fast decay of the coefficients in the polynomial expansion. The large coefficients usually occur in a lower multi-index set $I$, that is, if $\bm{\alpha}\in I$ and $\bm{\gamma} \le \bm{\alpha}$, then $\bm{\gamma}\in I$ \cite[Section 1.5]{adcock2022sparse}. To circumvent the curse of dimensionality,  the decay property of the expansion coefficients can be exploited and a hyperbolic cross sparse index set can be used to construct an approximation. The hyperbolic cross multi-index set $I$ with maximum order $p\in \mathbb{N}$ is defined by
\begin{equation} \label{eq:HC-defi}
    I := \Big\{\bm{\alpha} = (\alpha_j)_{j=1}^d \in \mathbb{N}_0^d: \prod_{j=1}^d \max(\alpha_j, 1) \le p\Big\},
\end{equation}
which has a cardinality $\mathcal{O}(p(\ln p)^{d-1})$. The best approximation error of the sparse Hermite polynomial approximation with hyperbolic cross index set was analyzed in \cite{luo2018error}; see \cref{sec:convergence} for details.

Now, we introduce a property of derivatives of Hermite polynomials, which plays an important role in reducing the computational cost. The first-order derivative of the one-dimensional normalized and generalized Hermite polynomial $H_n^{(t)}(y)$ satisfies
\begin{equation*}
    \frac{\mathrm{d}}{\mathrm{d}y}\left(H_n^{(t)}(y)\right) = \sqrt{\frac{n}{t}} H_{n-1}^{(t)}(y).
\end{equation*}
For the $d$-dimensional Hermite polynomial, we have
\begin{equation} \label{eq:derivative_herm}
    \frac{\partial}{\partial y_j}\left( H_{\bm{\alpha}}^{(t)}(\mathbf{y}) \right) = \sqrt{\frac{\alpha_j}{t}} H_{\bm{\alpha} - \mathbf{e}_j}^{(t)}(\mathbf{y}),
\end{equation}
where $\mathbf{e}_j$ is the $j$-th canonical basis vector. This implies that for a lower multi-index set $I \subset \mathbb{N}_0^d$, we have $\bm{\alpha} - \mathbf{e}_j \in I$ if $\bm{\alpha}\in I$. Thus, once the evaluations of $H_{\bm{\alpha}}^{(t)}(\mathbf{y})$ for $\bm{\alpha}\in I$ are available, the gradients of $H_{\bm{\alpha}}^{(t)}(\mathbf{y})$ for $\bm{\alpha}\in I$ can be evaluated cheaply.

\section{Algorithm and complexity} \label{sec:algorithm-cost}
Now we derive the main methodology to approximate the CVF $c_{t_k}$ by matching values of $u_{t_{k+1}}$, and analyze its computational complexity.
Below we abbreviate the notations $c_{t_k}$, $u_{t_k}$ and $\mathbf{W}_{t_k}$ to $c_k$, $u_k$ and $\mathbf{W}_k$, etc, for $k=0,1,\dots, N$.

\subsection{Gradient-enhanced Least Squares (G-LS)}
First we derive a linear backward stochastic differential equation (BSDE) for the CVF $ c_k$. It is worth noting that the corresponding BSDE does not contain a generator, and the option price is solely a function of Brownian motion rather than the underlying asset. This fact is crucial for developing the G-LSM using sparse Hermite polynomials.
\begin{theorem} \label{thm: linear bsde}
The CVF $c_k(\mathbf{W}_k)$ satisfies the linear BSDE
    \begin{equation*}
        c_k(\mathbf{W}_k) = u_{k+1}(\mathbf{W}_{k+1}) - \int_{t_k}^{t_{k+1}} \nabla_{\mathbf{w}}f(t, \mathbf{W}_t) \cdot \mathrm{d}\mathbf{W}_t,
    \end{equation*} 
where $f(t, \mathbf{w}), t\in [t_k, t_{k+1}]$, is defined by
$f(t, \mathbf{w}) := \mathbb{E}[ u_{k+1}(\mathbf{W}_{k+1}) | \mathbf{W}_t = \mathbf{w} ].$
\end{theorem}
\begin{proof}
This is a direct consequence of martingale representation theorem \cite[Theorem 2.5.2]{zhang2017backward}. For the sake of completeness, we provide a brief proof. 
By the Feynman-Kac formula, $f(t, \mathbf{w})$ satisfies a $d$-dimensional parabolic PDE subject to a terminal condition
\begin{equation}\label{eq:para-terminal}
\left\{
\begin{aligned}
\frac{\partial f}{\partial t}(t, \mathbf{w}) + \frac{1}{2}\operatorname{Tr}\left(\operatorname{Hess}_{\mathbf{w}}f(t,\mathbf{w}) \right) &= 0,\quad  t\in [t_k, t_{k+1}), \\
f(t_{k+1}, \mathbf{w}) &= u_{t_{k+1}}(\mathbf{w}).
\end{aligned}
\right.
\end{equation}
Let $Y_t := f(t, \mathbf{W}_t)$. By It\^{o}'s formula, we have
    \begin{equation*}
        \dd Y_t = \left( \frac{\partial f}{\partial t}(t, \mathbf{W}_t) + \frac{1}{2}\operatorname{Tr}\left(\operatorname{Hess}_{\mathbf{w}}f(t,\mathbf{W}_t) \right) \right) \,\mathrm{d}t + \nabla_{\mathbf{w}}f(t, \mathbf{W}_t) \cdot \mathrm{d}\mathbf{W}_t.
    \end{equation*}
In view of \cref{eq:para-terminal}, the drift term vanishes. After taking the stochastic integral and using the terminal condition in \cref{eq:para-terminal}, we obtain the desired assertion.
\end{proof}

\begin{remark}
The Feynman-Kac formula in the classical book \cite{oksendal2013stochastic} requires $C^2$ regularity of the terminal condition, which is not satisfied by the value function. The regularity requirement can be relaxed by considering the generalized result of nonlinear Feynman-Kac formula in \cite[p. 105]{zhang2017backward}, which only requires the terminal condition to be uniformly Lipschitz continuous.
\end{remark}

Next, we construct an approximation of the CVF ${c}_{k}$ by matching the terminal condition over the interval $[t_k,t_{k+1}]$. By \cref{thm: linear bsde}, the terminal value over $[t_k,t_{k+1}]$ can be approximated by the Euler discretization:
\begin{equation} \label{eq: approximation equation}
     \bar{u}_{k+1}(\mathbf{W}_{k+1}) =
     {c}^{{\rm CLS}}_{k}(\mathbf{W}_k) + \nabla {c}^{\op{CLS}}_{k}(\mathbf{W}_k) \cdot \Delta\mathbf{W}_k,
\end{equation}
where $\Delta\mathbf{W}_k:=\mathbf{W}_{k+1} - \mathbf{W}_k$ is the Brownian increment and ${c}^{\op{CLS}}_{k}$ denotes an approximation to the CVF ${c}_{k}$. Let $\hat{u}_{k+1}$ be the value function computed in the last time step. Then using $P_{I,k}$ defined in \cref{eq:pkI} and \cref{eq:HC-defi} as the ansatz space for ${c}^{\op{CLS}}_{k}$, we solve for ${c}^{\op{CLS}}_{k}$ using the least squares regression:
\begin{equation} \label{eq: expected risk}
    {c}_{k}^{\op{CLS}}(\mathbf{W}_k) = \argmin_{\psi \in P_{I,k}} E_k(\psi),
\end{equation}
with $E_k(\cdot): P_{I,k}\to \mathbb{R}^+$ being the quadratic loss (i.e., mean squared error) defined by
\begin{align*}
E_k(\psi):=\mathbb{E}\left[\left( \hat{u}_{k+1}(\mathbf{W}_{k+1}) - \psi(\mathbf{W}_k) - \nabla \psi(\mathbf{W}_k) \cdot \Delta \mathbf{W}_k \right)^2 \right].
\end{align*}
Finally, the numerical value $\hat{u}_k$ at time $t_k$ is updated to be the discounted exercise or continuation value at $t_k$:
\begin{align}\label{eq:num-val}
\hat{u}_{k}(\mathbf{W}_k)=\begin{cases}
    g_k(\mathbf{S}_k), \quad &\text{ if exercise}, \\
    c_k^{\op{CLS}}(\mathbf{W}_k), \quad &\text{ if continue}.
\end{cases}
\end{align}
Since the option is only profitable when exercised in the in-the-money region $\Omega_{\op{ITM}} = \{\mathbf{s}\in \R^d: g_k(\mathbf{s})>0\}$, we make the decision of exercising the option when $c_k^{\op{CLS}}(\mathbf{W}_k) < g_k(\mathbf{S}_k)$ and $\mathbf{S}_k \in \Omega_{\op{ITM}}$; otherwise, the option will be continued.

\subsection{Gradient-enhanced Least Squares Monte Carlo (G-LSM)} \label{subsec:glsm}
Now we derive the methodology based on the Monte Carlo method to solve \cref{eq: expected risk} numerically and analyze its computational complexity. Let $N_b=|I| $ and let $\{\phi_n^{k}(\mathbf{W}_k)\}_{n=1}^{N_b}$ be the set of Hermite polynomials in a scalar-indexed form. Then any function $\psi\in P_{I,k}$ can be expressed as
\begin{align*}
\psi(\mathbf{W}_k) = \sum_{n=1}^{N_b} \beta_n \phi_n^{k}(\mathbf{W}_k),
\end{align*}
with its gradient given by
\begin{equation} \label{eq: gradient of approx cv}
    \nabla \psi(\mathbf{W}_k) = \sum_{n=1}^{N_b} \beta_n \nabla \phi_n^{k}(\mathbf{W}_k).
\end{equation}
In practice, the continuous least squares problem \cref{eq: expected risk} is solved by minimizing its Monte Carlo approximation:
\begin{equation} \label{eq: empirical risk}
    \hat{c}_k := \argmin_{\psi \in P_{I,k}}
    \frac{1}{M}\sum_{m=1}^M
    \big(\hat{u}_{k+1}(\mathbf{W}^m_{k+1}) -
    \psi(\mathbf{W}^m_k) - \nabla \psi(\mathbf{W}^m_k) \cdot  \Delta\mathbf{W}^m_k
    \big)^2,
\end{equation}
where $\{\mathbf{W}^m_k\}_{m=1}^M$ are $M$ independent paths of the Gaussian random process $\mathbf{W}_k$ and $\Delta\mathbf{W}^m_k:=\mathbf{W}^m_{k+1} - \mathbf{W}^m_k$ is the $m$th path of the increment $\Delta \mathbf{W}_k$. Let $A_k\in\mathbb{R}^{M\times N_b}$, $\bm{\beta}_k\in\mathbb{R}^{ N_b}$ and $\mathbf{\hat{u}}_{k+1}\in\mathbb{R}^{ M}$ with their components defined by
\begin{equation} \label{eq:defi-linear-system}
    \begin{aligned}
        (A_k)_{m,n} &= \phi_n^{k}(\mathbf{W}^m_k) + \nabla \phi_n^{k}(\mathbf{W}^m_k) \cdot \Delta\mathbf{W}^m_k , \\
        (\bm{\beta}_k)_n &= \beta_n, \quad 
        (\mathbf{\hat{u}}_{k+1})_m = \hat{u}_{k+1}(\mathbf{W}^m_{k+1})
    \end{aligned}
\end{equation}
for $m = 1,\dots, M$, $n = 1,\dots, N_b$. 

Then finding the optimal polynomial in \cref{eq: empirical risk} amounts to solving the classical least squares problem
\begin{equation}\label{eq:ls-linearSystem}
    \bm{\beta}_k = \argmin_{\bm{\beta}} \|A_k \bm{\beta} - \mathbf{\hat{u}}_{k+1}\|_2^2 = (A_k^\top A_k)^{-1}A_k^\top \mathbf{\hat{u}}_{k+1}.
\end{equation}
The proposed algorithm is summarized in \cref{alg: sparse hermite polynomial expansions}.

\begin{remark}
Note that the well-known least squares Monte Carlo (LSM) method  \cite{longstaff2001valuing, tsitsiklis2001regression} also approximates the CVF with a finite number of basis functions. The orthonormal Hermite polynomials is one possible choice. However, LSM computes the coefficients by projecting the value, $u_{t_{k+1}}(\mathbf{W}_{t_{k+1}})$ in \cref{eq: continuation value W}, onto a finite-dimensional space spanned by the basis functions due to the projection nature of conditional expectation:
\begin{equation*}
    \hat{c}_k^{LSM} := \argmin_{\psi \in P_{I,k}}
    \frac{1}{M}\sum_{m=1}^M
    \big(\hat{u}_{k+1}(\mathbf{W}^m_{k+1}) -
    \psi(\mathbf{W}^m_k) \big)^2.
\end{equation*}
In contrast, G-LSM computes the coefficients by matching the approximate and exact value of $u_{t_{k+1}}(\mathbf{W}_{t_{k+1}})$, see \cref{eq: empirical risk}. 
\end{remark}

\begin{algorithm}[htbp]
\caption{G-LSM}
\label{alg: sparse hermite polynomial expansions}
\begin{algorithmic}[1]
\REQUIRE Market parameters: $ \mathbf{S}_0, r, \delta_i, \sigma_i, P $ \par
        Option parameters: payoff function $g(\mathbf{s})$ \par
        Algorithm parameters: $N, M, p$
\ENSURE Option price $v_0$
\STATE Compute the hyperbolic cross multi-index set $I$ with maximum order $p$ 
\STATE Generate $M$ sample paths
\STATE Initialize values $\mathbf{\hat{u}}_N = (e^{-rT}g(\mathbf{S}_N^m))_{m=1}^M$
\STATE Initialize stopping times $\tau_\star = T$
\FOR {$k = N-1:-1:1$}
    \STATE $\Phi_k = $ basis matrix$(\mathbf{W}_k, I)$ with $(\Phi_k)_{m,n} = \phi_n^k(\mathbf{W}_k^m)$
    \STATE Compute matrix $A_k$ with \cref{alg: compute matrix A}
    \STATE Solve system of linear equations: $A_k\bm{\beta}_k = \mathbf{u}_{k+1}$
    \STATE Update $(\mathbf{\hat{u}}_{k})_m = \begin{cases}
        e^{-rk\Delta t} g(\mathbf{S}^m_k), &\text{if exercise} \\
        (\Phi_k \bm{\beta}_k)_m, &\text{if continue}
    \end{cases}
    \text{ and }
    \tau_\star = \begin{cases}
        k\Delta t,  &\text{if exercise} \\
        \tau_\star, &\text{if continue}
    \end{cases}$
\ENDFOR
\STATE $v_0 = \max\left\{\frac{1}{M} \sum_{m=1}^M e^{-r\tau_\star^m} g(\mathbf{S}^m_{\tau_\star^m}),g(\mathbf{S}_0)\right\}$
\end{algorithmic}
\end{algorithm}

For the efficient computation of the matrix $A_k$, using \cref{eq:derivative_herm}, i.e.,
\begin{equation} \label{eq:basis-derivative}
    \frac{\partial \phi_n^{k}}{\partial w_j}(\mathbf{W}^m_k) = \sqrt{\frac{\alpha_j}{t_k}} \phi_{n'}^{k}(\mathbf{W}^m_k) \quad \text{ for } \phi_n^{k} = H_{\bm{\alpha}}^{(t_k)}, \phi_{n'}^{k} = H_{\bm{\alpha} - \mathbf{e}_j}^{(t_k)},
\end{equation}
we can compute the matrix $A_k$ from the basis matrix $\Phi_k$ and the increment $\Delta\mathbf{W}_k$. This routine is summarized in \cref{alg: compute matrix A}, where $(A_k)_{n}$ represents the $n$-th column of the matrix $A_k$.

Finally, we analyze the computational complexity of \cref{alg: sparse hermite polynomial expansions}. We employ backward induction \cref{eq: backward induction} with $N$ time steps to price American options, which has a complexity $\mathcal{O}(N)$. For each fixed $t_k$, the computation consists of steps 6-9 in \cref{alg: sparse hermite polynomial expansions}. Step 6 has a complexity $\mathcal{O}(MN_b)$, when using $M$ samples and $N_b$ basis functions. Step 7 is detailed in \cref{alg: compute matrix A}, which has a linear complexity in the number of nonzero $\alpha_j$ for all $\bm{\alpha} \in I$ and $j=1,\dots,d$, 
\begin{align*}
\|I\|_0:=\sum_{\bm{\alpha}\in I}\#\left\{j=1,\cdots,d: \alpha_j\neq 0\right\}. 
\end{align*}
Note that $\|I\|_0$ can be calculated by $(N_b - N_{b,p-1})d$ for maximum order $p$, where $N_{b,p-1}$ represents the number of basis functions with maximum order $p-1$. \Cref{fig: hc_nnz} shows that $\|I\|_0$ scales nearly linearly with respect to $N_b$. Thus, the complexity of \cref{alg: compute matrix A} is linear in $N_b$. Step 8 involves solving a linear system of size $M\times N_b$. Using an iterative solver, the cost is $\mathcal{O}(MN_b)$ if the matrix $A$ is well-conditioned. Step 9 involves matrix-vector multiplication, which has a complexity $\mathcal{O}(MN_b)$. Hence, with a fixed sample size $M$ and number $N$ of time steps, the total cost of the proposed G-LSM is $\mathcal{O}(NMN_b)$. 

\begin{algorithm}[htbp]
\caption{Compute the matrix $A_k$ using $I, \Phi_k$ and $\Delta\mathbf{W}_k$}
\label{alg: compute matrix A}
\begin{algorithmic}[1]
\STATE Initialize $A_k = \Phi_k$ 
\FOR{$j=1:d$}
    \FOR{$\bm{\alpha} \in I$}
        \IF{$\alpha_j \ge 1$}
        \STATE $(A_k)_{n} = (A_k)_{n} + (\Delta W_k)_j \odot (\Phi_k)_{n'} \sqrt{\frac{\alpha_j}{t_k}}$.
        \ENDIF
    \ENDFOR
\ENDFOR
\end{algorithmic}
\end{algorithm}

\begin{figure}[htbp]
    \centering
    \begin{tabular}{ccc}
    %\subfloat[ \label{fig: hc_nnz_order4}]{
    
    &%\subfloat[ \label{fig: hc_nnz_order6}]{
    
    &
    %\subfloat[ \label{fig: hc_nnz_order10}]{
    \\
    (a) $p=4$ & (b) $p=6$ & (c) $p=10$
    \end{tabular}
    \caption{$\|I\|_0$ scales linearly with respect to $N_b$.}
    \label{fig: hc_nnz}
\end{figure}


\subsection{Stability of the system of linear equations (\ref*{eq:ls-linearSystem})}
Next, we derive the stability of the system of linear equations \eqref{eq:ls-linearSystem} by estimating the condition number $\mathrm{cond}(A_k^\top A_k)$ (under the standard Euclidean norm) of the matrix $A_k^\top A_k$ defined in \eqref{eq:defi-linear-system} when the number $M$ of simulation paths is large. This result underpins the efficiency of the standard conjugate gradient method \cite{Hestenes:1952} for solving the least-squares problem (the associated normal equation, to be more precise).

\begin{theorem}\label{thm:cond}
Let $A_k\in\mathbb{R}^{M\times N_b}$ be defined in \eqref{eq:defi-linear-system}. Then there holds
\begin{align*}
\lim_{M\to\infty}\operatorname{cond}(A_k^\top A_k)= 1+\frac{p+d-1}{k} \quad\text{almost\; surely.}
\end{align*}
\end{theorem}
\begin{proof}
By definition, the $(i,j)$ entry of matrix $A_k^\top A_k$ is given by 
\begin{equation} \label{eq:entry-matrix}
    (A_k^\top A_k)_{i,j} = \sum_{m=1}^M \left( \phi_i^k(\mathbf{W}_k^m) + \nabla \phi_i^k(\mathbf{W}_k^m)\cdot \Delta \mathbf{W}_k^m\right) \left( \phi_j^k(\mathbf{W}_k^m) + \nabla \phi_j^k(\mathbf{W}_k^m)\cdot \Delta \mathbf{W}_k^m\right).
\end{equation}
By the vanilla Monte Carlo approximation, we obtain 
\begin{align}
        &\quad \lim_{M\to\infty} \frac{1}{M}\sum_{m=1}^M \left( \phi_i^k(\mathbf{W}_k^m) + \nabla \phi_i^k(\mathbf{W}_k^m)\cdot \Delta \mathbf{W}_k^m\right) \left( \phi_j^k(\mathbf{W}_k^m) + \nabla \phi_j^k(\mathbf{W}_k^m)\cdot \Delta \mathbf{W}_k^m\right)\nonumber \\
        &= \mathbb{E}\left[ \phi_i^k(\mathbf{W}_k)\phi_j^k(\mathbf{W}_k) + \nabla \phi_i^k(\mathbf{W}_k)\cdot \nabla\phi_j^k(\mathbf{W}_k) (\Delta \mathbf{W}_k)^2 \right]\nonumber \\
        &\quad + \mathbb{E}\left[ \left(\phi_i^k(\mathbf{W}_k) \nabla\phi_j^k(\mathbf{W}_k) + \phi_j^k(\mathbf{W}_k) \nabla\phi_i^k(\mathbf{W}_k) \right)\cdot \Delta \mathbf{W}_k \right] a.s. \nonumber \\
        &=  \mathbb{E}\left[ \phi_i^k(\mathbf{W}_k)\phi_j^k(\mathbf{W}_k) + \nabla \phi_i^k(\mathbf{W}_k)\cdot\nabla\phi_j^k(\mathbf{W}_k) (\Delta \mathbf{W}_k)^2 \right],\label{eq:entry-form}
\end{align}
where the last equality is due to the vanishing mean of Brownian motion increment $\mathbb{E}[\Delta \mathbf{W}_k] = 0$.

Furthermore, since $\phi_i^k$ and $\phi_j^k$ are generalized and normalized Hermite polynomials, then $\mathbb{E}[\phi_i^k(\mathbf{W}_k)\phi_j^k(\mathbf{W}_k)] = \delta_{ij}$. Then the identity $\mathbb{E}[(\Delta \mathbf{W}_k)^2] = \Delta t$ and \eqref{eq:basis-derivative} lead to 
\begin{equation*}
    \mathbb{E}[\nabla \phi_i^k(\mathbf{W}_k)\cdot\nabla\phi_j^k(\mathbf{W}_k) (\Delta \mathbf{W}_k)^2] = \bE\left[\sum_{\ell = 1}^d \partial_\ell \phi_i^k(\mathbf{W}_k)\partial_\ell \phi_j^k(\mathbf{W}_k)\right]\Delta t = 0,\quad \text{if }i\ne j.
\end{equation*}
This, together with \eqref{eq:entry-matrix} and \eqref{eq:entry-form}, implies that if $i\ne j$, then $\lim_{M\to\infty}(A_k^\top A_k)_{i,j} = 0$ a.s., i.e., the off-diagonal entries are nearly vanishing when $M\to \infty$.
To estimate the diagonal entries, let $\phi_i^k(\mathbf{W_k}) = H_{\bm{\alpha}}^{(t_k)}$ for some multi-index $\bm{\alpha}$. Using the first order derivative formula \eqref{eq:basis-derivative} and $t_k = k\Delta t$, we derive 
\begin{equation} \label{eq:element-AA}
    \begin{aligned}
        \bE\left[\sum_{\ell = 1}^d \partial_\ell \phi_i^k(\mathbf{W}_k)\partial_\ell \phi_i^k(\mathbf{W}_k)\right]\Delta t &= \Delta t\sum_{\ell = 1}^d \bE\left[ \left(\partial_\ell H_{\bm{\alpha}}^{(t_k)}(\mathbf{W}_k)\right)^2 \right] \\
        &=\Delta t\sum_{\ell = 1}^d \frac{\alpha_\ell}{t_k}\bE\left[ \left( H_{\bm{\alpha} - \mathbf{e}_\ell}^{(t_k)}(\mathbf{W}_k) \right)^2 \right] = \Delta t\frac{|\bm{\alpha}|}{t_k} = \frac{|\bm{\alpha}|}{k}.
    \end{aligned}
\end{equation}
By combining \eqref{eq:element-AA} with \eqref{eq:entry-form}, we obtain an estimate of the diagonal entry, \[\lim_{M\to\infty}\frac{1}{M}(A_k^\top A_k)_{i,i}= (1 + \frac{|\bm{\alpha}|}{k})\; a.s.
\] Here, the involved approximation error is only due to the Monte Carlo method. Thus, the matrix $A_k^\top A_k$ is approximately a diagonal matrix when the sample size $M$ is large. Hence, the definition of condition number implies
\begin{equation*}
     \lim_{M\to\infty}\operatorname{cond}(A_k^\top A_k)=\left(1 + \frac{\max_{\bm{\alpha \in I}}|\bm{\alpha}|}{k}\right) \Big/ \left(1 + \frac{\min_{\bm{\alpha \in I}}|\bm{\alpha}|}{k}\right) = 1 + \frac{p+d-1}{k} a.s.
\end{equation*}
Here, $I$ is the hyperbolic index set defined in \eqref{eq:HC-defi}. This proves the assertion.
\end{proof}

\begin{remark}\label{rmk:complexity}
We give a detailed comparison of the proposed G-LSM method with the numerical scheme in \cite{gobet2005regression} for solving BSDEs in terms of computational complexity. The scheme in \cite{gobet2005regression} solves the least squares problem 
    \begin{equation} \label{eq:ls-GLW05}
        \min_{\alpha_0, \alpha_1, \dots, \alpha_d} \frac{1}{M}\sum_{m=1}^M \left( Y_{t_{k+1}}^{N, I, I, M, m} - \alpha_0\cdot p_{0,k}^m - \sum_{l=1}^d \alpha_l \cdot p_{l,k}^m \Delta W_{l,k}^m \right)^2,
    \end{equation}
where each component of the BSDE solution $(Y_{t_k}, Z_{1, t_k}, \dots, Z_{d, t_k})$ is approximated by the linear combination of distinct basis functions, i.e., $\alpha_l\cdot p_{l,k}$. Suppose each approximation uses $N_b$ basis functions. Then $\alpha_l$ is a vector of length $N_b$ for each $l\in \{0, 1,\dots, d\}$. In total, the least squares problem \eqref{eq:ls-GLW05} involves $(d+1)\times N_b$ unknowns for a $d$-dimensional problem.  In contrast, the least squares problem in \eqref{eq: empirical risk} is equivalent to 
    \begin{equation} \label{eq:ls-glsm}
        \min_{\beta} \frac{1}{M}\sum_{m=1}^M \left( \hat{u}_{k+1}(\mathbf{W}_{k+1}^m) - \beta \cdot (A_k)_{m,\cdot} \right)^2,
    \end{equation}
    where only the CVF is approximated by the linear combination of sparse Hermite basis functions, i.e, the unknown $\beta \cdot \phi^k$ with $\beta\in \mathbb{R}^{N_b}$ and $\phi^k$ being the vector of basis functions. Here, $(A_k)_{m,\cdot}$ is the $m$-th row of the matrix $A_k$ in the linear system, which can be efficiently calculated by Algorithm 2. To the best of our knowledge, this idea is new in literature. We can see that the least squares problem \eqref{eq:ls-glsm} includes $N_b$ unknowns for a $d$-dimensional problem. Compared with $(d+1)\times N_b$ unknowns in the system \eqref{eq:ls-GLW05}, the least squares problem \eqref{eq:ls-glsm} incurs much lower computational complexity. Furthermore, the empirical regression matrix for the problem \eqref{eq:ls-GLW05} may not be invertible when the sample size becomes large as mentioned in \cite[page 2178]{gobet2005regression}, while the regression matrix for \eqref{eq:ls-glsm} is well-conditioned for large sample size as proved in Theorem \ref{thm:cond}.  In summary, in comparison with the scheme in \cite{gobet2005regression}, the proposed G-LSM method has $d$ times less computational complexity for  a $d$-dimensional problem, and  enjoys guaranteed stability when $M\to \infty$.
\end{remark}

\subsection{Computing deltas}
Now we present the approximation of deltas for constructing hedging strategies based on \cref{alg: sparse hermite polynomial expansions}. In the backward induction loop, \cref{eq: backward induction} and \cref{eq: defi greeks} imply the deltas at time $t_k$ is given by
\begin{equation*}
    \frac{\partial v_k(\mathbf{s})}{\partial s_j} = \left\{
        \begin{aligned}
            \frac{\partial g_k(\mathbf{s})}{\partial s_j}, \text{ if exercise}, \\
            \frac{\partial {\tilde{c}}_k(\mathbf{s})}{\partial s_j}, \text{ if continue}, 
        \end{aligned}
    \right.\quad j = 1,\dots, d.
\end{equation*}
In the continuation region, $\tilde{c}_k(\mathbf{s}) = c_k(\mathbf{w})$ with $\mathbf{s} = \mathbf{s}_0 \odot \exp(Q(\mu t_k + \sqrt{\Lambda}\mathbf{w}))$.
Once we have obtained the coefficients $\bm{\beta}_k$ of the expansion, the gradient $\nabla c_k(\mathbf{w})$ of the continuation value $c_k(\mathbf{w})$ is approximated by $\nabla \hat{c}_k(\mathbf{w})$ in \cref{eq: gradient of approx cv}. Hence, we calculate the deltas in the continuation region by
\begin{equation*}
    \frac{\partial {\tilde{c}}_k(\mathbf{s})}{\partial s_j} \approx \sum_{i=1}^d \frac{\partial \hat{c}_k(\mathbf{w})}{\partial w_i}\frac{\partial w_i}{\partial s_j}, \quad k = 1,\dots, {N-1}.
\end{equation*}
The gammas $\Gamma_{t_k}^{ij}$ can be approximated in a similar way.

Particularly, the Greeks at time $t=0$ involves solving an additional linear system, which arises from minimizing the $L_{\omega_{t_1}}^2(\R^d)$ error of approximating $u_{1}(\mathbf{W}_1)$. We employ a slightly different ansatz from the case of $t_k>0$, since the expansions in $H_{\bm{\alpha}}^{(0)}$ is not applicable. Instead, we take 
\begin{equation*}
    \hat{c}_0(\mathbf{w}) = \sum_{\bm{\alpha}\in I} \beta_{\bm{\alpha}} H_{\bm{\alpha}}^{(\Delta t)}(\mathbf{w}).
\end{equation*}
Since $c_0(\mathbf{W}_0)$ is a deterministic value given $\mathbf{W}_0 = \mathbf{0}$, we consider the Euler approximation
\begin{equation*}
    u_1(\mathbf{W}_1) \approx \hat{c}_0(\mathbf{0}) + \nabla \hat{c}_0(\mathbf{0}) \cdot \mathbf{W}_1.
\end{equation*}
Similar to \cref{eq: empirical risk}, the coefficients $\beta_{\bm{\alpha}}$ can be approximated by solving an $M\times N_b$ linear system. After that, the delta at $t=0$ is approximated by
\begin{equation*}
    \frac{\partial {\tilde{c}}_0(\mathbf{s}_0)}{\partial s_j} \approx \sum_{i=1}^d \frac{\partial \hat{c}_0(\mathbf{0})}{\partial w_i}\frac{\partial w_i}{\partial s_j}(\mathbf{s}_0),\quad \mbox{with }\mathbf{s} = \mathbf{s}_0 \odot  \exp(Q\sqrt{\Lambda}\mathbf{w}).
\end{equation*}

\section{Convergence analysis} \label{sec:convergence}
Now we analyze the convergence of \cref{alg: sparse hermite polynomial expansions}. 
We restrict the theoretical analysis to the multi-asset Black-Scholes model presented in \cref{sec:multi model} since the corresponding CVF can be reformulated as a function of independent Brownian motions and the associated best approximation error in the sparse Hermite polynomial ansatz space $P_{I,k}$ can be rigorously established, cf.  \cref{lem:best-error-HC}. In practice, \cref{alg: sparse hermite polynomial expansions} can be applied to general models driven by It\^{o} diffusions.  
We assume the Lipschitz continuity of the discounted payoff function $g_k(\cdot)$.
\begin{assumption}\label{assu:lip}
    The discounted payoff function $g_k(\cdot)$ is $L$-Lipschitz continuous: for  $k = 0,1,\dots,N$
    \begin{equation} \label{eq: Lipschitz assum}
        |g_k(\mathbf{s}) - g_k(\mathbf{s}')| \le L\|\mathbf{s} - \mathbf{s}'\|, \quad \forall \mathbf{s}, \mathbf{s}' \in \R^d.
    \end{equation}
\end{assumption}
Recall that $\hat{u}_{k}(\mathbf{w})$ is the numerical value function at time $t_k$ computed by \cref{alg: sparse hermite polynomial expansions}, which approximates the exact value ${u}_{k}(\mathbf{w})$ for $k = 0,1,\dots,N-1$. We aim to establish an upper bound for
$\max_{0\le k\le N-1} \|u_k -\hat{u}_k\|_{{L}^2_{\omega_k}(\R^d)}^2$.

\subsection{One-step error estimation}
First, we analyze the one step error over the interval $[t_k,t_{k+1}]$ for $k = 0,1,\dots,N-1$, i.e., the numerical value function $\hat{u}_{k}(\mathbf{w})$ as an approximation to the exact one ${u}_{k}(\mathbf{w})$. Given the previous value function $\hat{u}_{k+1}$, we define the current CVF  $ \check{c}_k(\mathbf{W}_k)$ by
\begin{equation}\label{eq:apprx-conti}
    \check{c}_k(\mathbf{W}_k) = \mathbb{E}[\hat{u}_{k+1}(\mathbf{W}_{k+1}) | \mathbf{W}_k].
\end{equation}
The classic backward Euler scheme for BSDE approximates $\nabla\check{c}_k$ by
\begin{equation}\label{eq:grad-est}
     \check{Z}_k:=\frac{1}{\Delta t} \mathbb{E}[\hat{u}_{k+1}(\mathbf{W}_{k+1})\Delta \mathbf{W}_k | \mathbf{W}_k].
\end{equation}
With the sparse Hermite polynomial ansatz space $P_{I,k}$, let $c_k^* \in P_{I,k}$ be the best approximation to $\check{c}_k$ in $H^1_{\omega_k}(\mathbb{R}^d)$ defined by
\begin{equation}\label{eq:best-approx}
    c^*_k := \argmin_{\psi \in P_{I,k}} \|\check{c}_k - \psi\|_{H^1_{\omega_k}(\R^d)}.
\end{equation}
Also we define the best approximation error $\mathcal{E}_k^{\rm best}$ by
\begin{equation} \label{eq: defi best approx}
    \mathcal{E}_k^{\rm best} := \|\check{c}_k - c_k^*\|_{H^1_{\omega_k}(\R^d)}^2.
\end{equation}
Next, we denote the statistical error of solving the discrete least squares problem \cref{eq: empirical risk} by
\begin{align}\label{eq:stat-lsm}
\mathcal{E}^{\rm stat}_k := \|\hat{c}_k - c_k^{\op{CLS}}\|_{{L}^2_{\omega_k}(\R^d)}^2,
\end{align}
where $c_k^{\rm CLS}$ solves the continuous least squares problem \cref{eq: expected risk} for $k = 0,1,\dots,N-1$.
In \cref{thm: one-step error} below, we derive the convergence of $c_k^{\op{CLS}}$ to the exact $c_k$ provided that the previous approximation error $\|u_{k+1} - \hat{u}_{k+1}\|_{L_{\omega_{k+1}}^2(\R^d)}^2$,  $\mathcal{E}_k^{\rm best}$ in \cref{eq: defi best approx}, and time step size $\Delta t$ are small. We first provide in \cref{lem:best-error-HC} an estimate of $\mathcal{E}_k^{\rm best}$ in \cref{eq: defi best approx}, which is a direct application of \cite[Theorem 4.2]{luo2018error}.

\begin{lemma} \label{lem:best-error-HC}
    For the hyperbolic cross index set $I$ with maximum order $p$ defined in \cref{eq:HC-defi}, given $\check{c}_k \in \mathcal{K}^m_{\omega_k}(\R^d)$, we have
    \begin{equation*}
        \mathcal{E}_k^{\rm best} \le C(m, d)\frac{1}{p^{m-1}} |\check{c}_k|_{\mathcal{K}^m_{\omega_k}(\R^d)}^2,
    \end{equation*}
    where $\mathcal{K}^m_{\omega_k}(\R^d)$ is the weighted Koborov-type space defined by
    \begin{equation*}
        \mathcal{K}^m_{\omega_k}(\R^d) = \{u: \partial^{\bm{\alpha}} u \in {L}^2_{\omega_k}(\R^d), 0\le |\bm{\alpha}|_\infty \le m\},
    \end{equation*}
    and $|\cdot|_{\mathcal{K}^m_{\omega_k}(\R^d)}$ is the seminorm defined by
    \begin{equation*}
        |u|_{\mathcal{K}^m_{\omega_k}(\R^d)} = \left( \sum_{|\bm{\alpha}|_\infty = m} \|\partial^{\bm{\alpha}} u\|_{{L}^2_{\omega_k}(\R^d)}^2 \right)^{1/2}.
    \end{equation*}
\end{lemma}

The next result provides a one-step error estimate of the $Z$ component in the BSDE by the backward Euler scheme \cref{eq:grad-est}. The proof is inspired by the idea of \cite[pages 816-817]{gobet2007error}.
By martingale representation theorem, the definition of $\check{c}_k$ \cref{eq:apprx-conti} implies the existence of a square integrable process $\tilde{Z}_s$, $t_k\le s\le t_{k+1}$, such that
\begin{equation} \label{eq: martingale representation}
    \hat{u}_{k+1}(\mathbf{W}_{k+1}) = \check{c}_k(\mathbf{W}_k) + \int_{t_k}^{t_{k+1}} \tilde{Z}_t \cdot \dd \mathbf{W}_t.
\end{equation}
\begin{lemma} \label{lem:one-step Z}
Let $\check{Z}_k$ and $\tilde{Z}_t$ be defined in \cref{eq:grad-est} and \cref{eq: martingale representation}. Then for some constant $C_1>0$, \begin{equation*}
        \mathbb{E}\left[ (\check{Z}_k - \tilde{Z}_{t_k})^2 \right] \le C_1 \Delta t^2.
    \end{equation*}
\end{lemma}
\begin{proof}
By the integration by parts formula of Malliavin calculus, we obtain
    \begin{equation*}
        \check{Z}_k = \frac{1}{\Delta t} \mathbb{E}\left[ \hat{u}_{k+1}(\mathbf{W}_{k+1}) \Delta \mathbf{W}_k | \mathbf{W}_k \right] = \frac{1}{\Delta t} \mathbb{E}\left[ \int_{t_k}^{t_{k+1}} D_t\mathbf{W}_{k+1} \nabla \hat{u}_{k+1}(\mathbf{W}_{k+1}) \,\dd t \bigg| \mathbf{W}_k \right],
    \end{equation*}
    where $D_t \mathbf{W}_{k+1}$ is an identity matrix of size $d\times d$.
Together with the identity
    \begin{equation*}
        \tilde{Z}_{t_k} = \nabla \check{c}_k(\mathbf{W}_k) = \frac{1}{\Delta t}\int_{t_k}^{t_{k+1}} \nabla \check{c}_k(\mathbf{W}_k) \,\dd t,
    \end{equation*}
it further leads to
    \begin{align*}
        \check{Z}_k - \tilde{Z}_{t_k} &= \frac{1}{\Delta t}\mathbb{E}\left[ \int_{t_k}^{t_{k+1}} \nabla \hat{u}_{k+1}(\mathbf{W}_{k+1}) - \nabla \check{c}_k(\mathbf{W}_k) \,\dd s \bigg| \mathbf{W}_k \right].
    \end{align*}
Then, by \cite[Equation (26)]{gobet2007error}, we obtain
    \begin{equation*}
        \check{Z}_k - \tilde{Z}_{t_k} = \int_{t_k}^{t_{k+1}} \mathbb{E}\left[ G(t, \mathbf{W}_t) |\mathbf{W}_k \right] \,\dd t,
    \end{equation*}
    for a bounded function $G$. Hence, there holds $|\check{Z}_k - \tilde{Z}_{t_k}| = \mathcal{O}(\Delta t)$, and there exists a constant $C_1>0$ such that $\mathbb{E}\left[ (\check{Z}_k - \tilde{Z}_{t_k})^2) \right] \le C_1 \Delta t^2$.
\end{proof}

Next, we provide an error estimate of solving the continuous least squares problem \cref{eq: expected risk} in terms of the error of the previous value function, the best approximation error in the sparse Hermite polynomial ansatz space \cref{eq: defi best approx} and the time step size $\Delta t$. The proof is inspired by the foundational work \cite{hure2020deep}. We first derive a one-step error propagation by using Young's inequality and then apply discrete Gronwall's inequality to establish global convergence. In contrast to the approach in \cite{hure2020deep}, our analysis also incorporates the statistical error introduced by the Monte Carlo method and the one-step error proven in \cref{lem:one-step Z}. Moreover, the best approximation error of sparse Hermite polynomial approximation is given in \cref{lem:best-error-HC}.
\begin{theorem} \label{thm: one-step error}
For $\Delta t \le 1/4$, with the constant $C_1>0$ in \cref{lem:one-step Z}, there holds
    \begin{equation*}
        \|c_k - c_k^{\op{CLS}}\|_{{L}^2_{\omega_k}(\R^d)}^2 \le (1 + 4 \Delta t) \|u_{k+1} - \hat{u}_{k+1}\|_{{L}^2_{\omega_{k+1}}(\R^d)}^2 + \tfrac{1}{2\Delta t} \mathcal{E}_k^{\rm best} + C_1 \Delta t^2.
    \end{equation*}
\end{theorem}
\begin{proof}
By the inequality $(a+b)^2 \le (1+ 4 \Delta t) a^2 + (1+\frac{1}{4 \Delta t})b^2$, we have
\begin{align*}
    \|c_k - c_k^{\op{CLS}}\|_{{L}^2_{\omega_k}(\R^d)}^2 &\le (1+4 \Delta t) \|c_k - \check{c}_k\|_{{L}^2_{\omega_k}(\R^d)}^2 + (1+\tfrac{1}{4\Delta t}) \|\check{c}_k - c_k^{\rm CLS}\|_{{L}^2_{\omega_k}(\R^d)}^2 \\
    &\le (1+4 \Delta t) \|c_k - \check{c}_k\|_{{L}^2_{\omega_k}(\R^d)}^2 + \tfrac{1}{2\Delta t} \|\check{c}_k - c_k^{\op{CLS}}\|_{{L}^2_{\omega_k}(\R^d)}^2,
\end{align*}
since $\Delta t \le 1/4$. For the first term, the definitions \cref{eq: continuation value W} and \cref{eq:apprx-conti}, the tower property, and Jensen's inequality lead to
\begin{equation*}
\begin{aligned}
    \mathbb{E}\left[ \left( c_k(\mathbf{W}_k) - \check{c}_k(\mathbf{W}_k) \right)^2 \right] &= \mathbb{E}\left[ \left(\mathbb{E}\left[ u_{k+1}(\mathbf{W}_{k+1}) - \hat{u}_{k+1}(\mathbf{W}_{k+1})|\mathbf{W}_k \right] \right)^2 \right] \\
    &\le \mathbb{E}\left[ \left(u_{k+1}(\mathbf{W}_{k+1}) - \hat{u}_{k+1}(\mathbf{W}_{k+1}) \right)^2 \right] = \|u_{k+1} - \hat{u}_{k+1}\|_{{L}^2_{\omega_{k+1}}(\R^d)}^2.
\end{aligned}
\end{equation*}
Hence, it remains to show
\begin{equation} \label{eq:remains-to-show}
    \|\check{c}_k - c_k^{\op{CLS}}\|_{{L}^2_{ \omega_{k}}(\R^d)}^2 \le \mathcal{E}_k^{\rm best} + 2C_1\Delta t^3.
\end{equation}
By plugging \cref{eq: martingale representation} into the quadratic loss function $E_k(\cdot)$ in \cref{eq: expected risk}, we obtain
\begin{align*}
    E_k(\psi) &= \mathbb{E}\left[ \left( \check{c}_k(\mathbf{W}_k) - \psi(\mathbf{W}_k) + \int_{t_k}^{t_{k+1}} \tilde{Z}_s\cdot \dd \mathbf{W}(s) - \nabla \psi(\mathbf{W}_k)\cdot \Delta \mathbf{W}_k \right)^2 \right] \\
    &= \mathbb{E}\left[ \left( \check{c}_k(\mathbf{W}_k) - \psi(\mathbf{W}_k) \right)^2 \right] + \mathbb{E}\left[ \left( \int_{t_k}^{t_{k+1}} \tilde{Z}_s\cdot \dd \mathbf{W}(s) - \nabla \psi(\mathbf{W}_k)\cdot \Delta \mathbf{W}_k \right)^2 \right].
\end{align*}
Now It\^{o} isometry implies the identity
\begin{align*}
    & \mathbb{E}\left[ \left( \int_{t_k}^{t_{k+1}} \tilde{Z}_s\cdot \dd \mathbf{W}(s) - \nabla \psi(\mathbf{W}_k)\cdot \Delta \mathbf{W}_k \right)^2 \right] \\
    =& \mathbb{E}\left[ \left( \int_{t_k}^{t_{k+1}} \tilde{Z}_s\cdot \dd \mathbf{W}(s) - \int_{t_k}^{t_{k+1}} \check{Z}_k \cdot \dd \mathbf{W}(s) + \check{Z}_k\cdot \Delta \mathbf{W}_k - \nabla \psi(\mathbf{W}_k)\cdot \Delta \mathbf{W}_k \right)^2 \right]\\
    =& \mathbb{E}\left[ \int_{t_k}^{t_{k+1}} |\tilde{Z}_s - \check{Z}_k|^2\,\dd s \right] + \Delta t \mathbb{E}\left[ \left|\check{Z}_k - \nabla \psi(\mathbf{W}_k) \right|^2 \right] \\
    &+ 2\mathbb{E}\left[ \int_{t_k}^{t_{k+1}} (\tilde{Z}_s - \check{Z}_k)\,\dd s \right] \cdot \mathbb{E}[\check{Z}_k - \nabla \psi(\mathbf{W}_k)].
\end{align*}
The definitions of $\check{Z}_k$ and $\tilde{Z}_s$ in \cref{eq:grad-est} and \cref{eq: martingale representation}, together with It\^{o} isometry, yield
\begin{align*}
\check{Z}_k = \frac{1}{\Delta t} \mathbb{E}\left[ \int_{t_k}^{t_{k+1}} \tilde{Z}_s \,\dd s \bigg| \mathbf{W}_k \right],
\end{align*}
which implies
\begin{align*}
\mathbb{E}\left[ \int_{t_k}^{t_{k+1}} (\tilde{Z}_s - \check{Z}_k)\,\dd s \right] = 0.
\end{align*}
Hence, we can express the loss function $E_k(\psi)$ as
\begin{equation*}
    E_k(\psi) = \mathbb{E}\left[ \left( \check{c}_k(\mathbf{W}_k) - \psi(\mathbf{W}_k) \right)^2 \right] + \mathbb{E}\left[ \int_{t_k}^{t_{k+1}} |\tilde{Z}_s - \check{Z}_k|^2\,\dd s \right] + \Delta t \mathbb{E}\left[ \left|\check{Z}_k - \nabla \psi(\mathbf{W}_k) \right|^2 \right].
\end{equation*}
The equality \cref{eq: expected risk} implies that $E_k(c_k^{\rm CLS}) \le E_k(\psi)$ for all $\psi\in P_{I,k}$. By taking $\psi := c_k^*$ as defined in \cref{eq:best-approx}, we obtain
\begin{align*}
    &\quad \mathbb{E}\left[ \left( \check{c}_k(\mathbf{W}_k) - c_k^{\rm CLS}(\mathbf{W}_k) \right)^2 \right] + \Delta t \mathbb{E}\left[ \left(\check{Z}_k - \nabla c_k^{\rm CLS}(\mathbf{W}_k) \right)^2 \right] \\
    &\le \mathbb{E}\left[ \left( \check{c}_k(\mathbf{W}_k) - c_k^*(\mathbf{W}_k) \right)^2 \right] + \Delta t \mathbb{E}\left[ \left(\check{Z}_k - \nabla c_k^*(\mathbf{W}_k) \right)^2 \right].
\end{align*}
Note that $\tilde{Z}_{t_k} = \nabla \check{c}_k(\mathbf{W}_k)$. By subtracting and adding $\tilde{Z}_{t_k}$, we have
\begin{align*}
    &\quad \|\check{c}_k - c_k^{\op{CLS}}\|_{{L}^2_{\omega_{k}}(\R^d)}^2 \\
    &\le \|\check{c}_k - c_k^*\|_{{L}^2_{\omega_{k}}(\R^d)}^2 + 2\Delta t \mathbb{E}\left[ (\check{Z}_k - \tilde{Z}_{t_k})^2 \right] + 2\Delta t \mathbb{E}\left[ |\nabla \check{c}_k(\mathbf{W}_k) - \nabla c_k^*(\mathbf{W}_k)|^2 \right] \\
    &\le \|\check{c}_k - c_k^*\|_{H^1_{ \omega_{k}}(\R^d)}^2 + 2C_1\Delta t^3,
\end{align*}
where the last inequality follows from the condition $2\Delta t \le 1$ and \cref{lem:one-step Z}. This shows the estimate \cref{eq:remains-to-show}, and completes the proof of the theorem.
\end{proof}

Now, we can give the one-step error propagation of $\|u_k - \hat{u}_k\|_{{L}^2_{\omega_k}(\R^d)}^2$ given $\|u_{k+1} - \hat{u}_{k+1}\|_{{L}^2_{\omega_{k+1}}(\R^d)}^2$.
\begin{corollary} \label{cor: one-step error}
For $\Delta t\le 1/6$, there holds
    \begin{equation*}
        \|u_k - \hat{u}_k\|_{{L}^2_{\omega_k}(\R^d)}^2 \le (1+6 \Delta t) \|u_{k+1} - \hat{u}_{k+1}\|_{{L}^2_{\omega_{k+1}}(\R^d)}^2 + \left(1 + \frac{1}{2\Delta t}\right) \mathcal{E}_k^{\rm best} + \frac{6}{5}C_1\Delta t^2 + \frac{6}{5\Delta t} \mathcal{E}_k^{\rm stat},
    \end{equation*}
    where $\mathcal{E}_k^{\rm best}$ and $\mathcal{E}_k^{\rm stat}$ are defined in \cref{eq: defi best approx} and \cref{eq:stat-lsm}, and $C_1>0$ is the constant in \cref{lem:one-step Z}.
\end{corollary}
\begin{proof}
First, using inequality $(a+b)^2 \ge (1-\Delta t)a^2 - \frac{1}{\Delta t}b^2$ for any $a,b\in\mathbb{R}$, leads to
\begin{equation*}
    \|c_k - c_k^{\op{CLS}}\|_{{L}^2_{\omega_k}(\R^d)}^2 \ge (1-\Delta t) \|c_k - \hat{c}_k\|_{{L}^2_{\omega_k}(\R^d)}^2 - \frac{1}{\Delta t}\|\hat{c}_k - c_k^{CLS}\|_{{L}^2_{\omega_k}(\R^d)}^2.
\end{equation*}
This and \cref{thm: one-step error} yield
\begin{align*}
    &\quad \|c_k - \hat{c}_k\|_{{L}^2_{\omega_k}(\R^d)}^2 \\
    &\le \frac{1+4\Delta t}{1-\Delta t} \|u_{k+1} - \hat{u}_{k+1}\|_{{L}^2_{\omega_{k+1}}(\R^d)}^2 + \frac{1}{2\Delta t(1-\Delta t)} \mathcal{E}_k^{\rm best} + \frac{C_1\Delta t^2}{1-\Delta t} + \frac{1}{\Delta t(1-\Delta t)} \mathcal{E}_k^{\rm stat} \\
    &\le (1 + 6\Delta t) \|u_{k+1} - \hat{u}_{k+1}\|_{{L}^2_{\omega_{k+1}}(\R^d)}^2 + \left(1 + \frac{1}{2\Delta t} \right) \mathcal{E}_k^{\rm best} + \frac{6}{5}C_1\Delta t^2 + \frac{6}{5\Delta t} \mathcal{E}_k^{\rm stat},
\end{align*}
provided that $\Delta t\le 1/6$.
Finally, using inequality $|\max(a,b) - \max(a,c)| \le |b-c|$, we derive
\begin{equation*}
    \|u_k - \hat{u}_k\|_{{L}^2_{\omega_k}(\R^d)}^2 = \|\max(G_k, c_k) - \max(G_k, \hat{c}_k)\|_{{L}^2_{\omega_k}(\R^d)}^2 \le \|c_k - \hat{c}_k\|_{{L}^2_{\omega_k}(\R^d)}^2.
\end{equation*}
This completes the proof.
\end{proof}

\subsection{Global error estimation}
Finally, we prove a global error estimate.

\begin{theorem} \label{thm: global error}
For $\Delta t \le 1/6$, there holds
    \begin{equation*}
        \max_{0\le k\le N-1} \|u_k - \hat{u}_k\|_{{L}^2_{\omega_k}(\R^d)}^2 \le C\left( \Delta t +  N\sum_{k=0}^{N-1} \left( \mathcal{E}_k^{\rm best} + \mathcal{E}_k^{\rm stat} \right) \right),
    \end{equation*}
    where the constant $C = \max\left(\frac{6}{5T}e^{6T}, \frac{6T}{5}C_1 e^{6T} \right)$ only depends on the finite time horizon $T$ and the constant $C_1>0$ defined in \cref{lem:one-step Z}.
\end{theorem}

\begin{proof}
    For $\Delta t\le 1/6$, \cref{cor: one-step error} implies
    \begin{equation*}
        \|u_k - \hat{u}_k\|_{{L}^2_{\omega_k}(\R^d)}^2 \le (1+6 \Delta t) \|u_{k+1} - \hat{u}_{k+1}\|_{{L}^2_{\omega_{k+1}}(\R^d)}^2 + \frac{6}{5\Delta t} \left( \mathcal{E}_k^{\rm best} + \mathcal{E}_k^{\rm stat} \right) + \frac{6}{5}C_1\Delta t^2.
    \end{equation*}
    Using the discrete Gronwall's inequality and $\Delta t = T/N$, we obtain
    \begin{equation*}
        \max_{0\le k\le N-1} \|u_k - \hat{u}_k\|_{{L}^2_{\omega_k}(\R^d)}^2 \le e^{6T} \|u_N - \hat{u}_N\|_{{L}^2_{\omega_N}(\R^d)}^2 + e^{6T} \frac{6N}{5T} \sum_{k=0}^{N-1} \left( \mathcal{E}_k^{\rm best} + \mathcal{E}_k^{\rm stat} \right) + e^{6T} \frac{6T}{5}C_1 \Delta t.
    \end{equation*}
    Since the terminal condition $u_N$ is known, the result follows by $\|u_N - \hat{u}_N\|_{{L}_{\omega_N}^2(\R^d)} = 0$.
\end{proof}

\Cref{thm: global error} implies that the numerical value function $\hat{u}_{k}$ computed by \cref{alg: sparse hermite polynomial expansions} at each time $t_k$ approximates the exact one ${u}_{k}$ well if $\Delta t$, $\mathcal{E}_k^{\rm best}$, and $\mathcal{E}_k^{\rm stat}$ are small. Using \cref{lem:best-error-HC}, the best approximation error $\mathcal{E}_k^{\rm best}$ decays geometrically with the increase of the order $p$ of sparse Hermite polynomials. By the law of large numbers, the statistical error
$\mathcal{E}_k^{\rm stat}$ between the discrete and the continuous least squares will be small when using a large number of sample paths.

\begin{thebibliography}{99}
\bibitem[1]{adcock2022sparse} {\sc B.~Adcock, S.~Brugiapaglia, and C.~G. Webster}, {\em {Sparse Polynomial Approximation of High-Dimensional Functions}}, SIAM, Philadelphia, PA, 2022.
\bibitem[2]{bayer2023pricing} {\sc C.~Bayer, M.~Eigel, L.~Sallandt, and P.~Trunschke}, {\em Pricing high-dimensional {B}ermudan options with hierarchical tensor formats}, SIAM Journal on Financial Mathematics, 14 (2023), pp.~383--406.
\bibitem[3]{becker2019deep} {\sc S.~Becker, P.~Cheridito, and A.~Jentzen}, {\em Deep optimal stopping}, The Journal of Machine Learning Research, 20 (2019), pp.~2712--2736.
\bibitem[4]{becker2020pricing} {\sc S.~Becker, P.~Cheridito, and A.~Jentzen}, {\em Pricing and hedging {A}merican-style options with deep learning}, Journal of Risk and Financial Management, 13 (2020), p.~158.
\bibitem[5]{bouchard2012monte} {\sc B.~Bouchard and X.~Warin}, {\em {M}onte-{C}arlo valuation of {A}merican options: facts and new algorithms to improve existing methods}, in Numerical Methods in Finance: Bordeaux, June 2010, Springer, Berlin, 2012, pp.~215--255.
\bibitem[6]{chen2021deep} {\sc Y.~Chen and J.~W. Wan}, {\em Deep neural network framework based on backward stochastic differential equations for pricing and hedging {A}merican options in high dimensions}, Quantitative Finance, 21 (2021), pp.~45--67.
\bibitem[7]{E2017deepbsde} {\sc W.~E, J.~Han, and A.~Jentzen}, {\em Deep learning-based numerical methods for high-dimensional parabolic partial differential equations and backward stochastic differential equations}, Communications in Mathematics and Statistics, 5 (2017), pp.~349--380.
\bibitem[8]{el1997reflected} {\sc N.~El~Karoui, C.~Kapoudjian, E.~Pardoux, S.~Peng, and M.-C. Quenez}, {\em Reflected solutions of backward {SDE}'s, and related obstacle problems for {PDE}'s}, The Annals of Probability, 25 (1997), pp.~702--737.
\bibitem[10]{gao2023convergence} {\sc C.~Gao, S.~Gao, R.~Hu, and Z.~Zhu}, {\em Convergence of the backward deep bsde method with applications to optimal stopping problems}, SIAM Journal on Financial Mathematics, 14 (2023), pp.~1290--1303.
\bibitem[11]{gobet2007error} {\sc E.~Gobet and C.~Labart}, {\em Error expansion for the discretization of backward stochastic differential equations}, Stochastic Processes and their Applications, 117 (2007), pp.~803--829.
\bibitem[12]{gobet2005regression} {\sc E.~Gobet, J.-P. Lemor, and X.~Warin}, {\em A regression-based {M}onte {C}arlo method to solve backward stochastic differential equations}, The Annals of Applied Probability, (2005).
\bibitem[13]{guyon2013nonlinear} {\sc J.~Guyon and P.~Henry-Labordère}, {\em Nonlinear option pricing}, Chapman \& Hall/CRC financial mathematics series, CRC Press, Boca Raton, 2014 - 2014.
\bibitem[14]{Hestenes:1952} {\sc M.~R. Hestenes and E.~Stiefel}, {\em Methods of conjugate gradients for solving linear systems}, J. Research Nat. Bur. Standards, 49 (1952), pp.~409--436.
\bibitem[15]{hure2020deep} {\sc C.~Hur{\'e}, H.~Pham, and X.~Warin}, {\em Deep backward schemes for high-dimensional nonlinear {PDE}s}, Mathematics of Computation, 89 (2020), pp.~1547--1579.
\bibitem[17]{lapeyre2021neural} {\sc B.~Lapeyre and J.~Lelong}, {\em Neural network regression for {B}ermudan option pricing}, Monte Carlo Methods and Applications, 27 (2021), pp.~227--247.
\bibitem[18]{longstaff2001valuing} {\sc F.~Longstaff and E.~Schwartz}, {\em Valuing {A}merican options by simulation: a simple least-squares approach}, The Review of Financial Studies, 14 (2001), pp.~113--147.
\bibitem[19]{ludkovski2018kriging} {\sc M.~Ludkovski}, {\em Kriging metamodels and experimental design for {B}ermudan option pricing}, Journal of Computational Finance, 22 (2018), pp.~37--77.
\bibitem[20]{luo2018error} {\sc X.~Luo}, {\em Error analysis of the {W}iener--{A}skey polynomial chaos with hyperbolic cross approximation and its application to differential equations with random input}, Journal of Computational and Applied Mathematics, 335 (2018), pp.~242--269.
\bibitem[23]{oksendal2013stochastic} {\sc B.~Oksendal}, {\em Stochastic differential equations: an introduction with applications}, Springer Science \& Business Media, 2013.
\bibitem[26]{seydel2006tools} {\sc R.~Seydel and R.~Seydel}, {\em Tools for computational finance}, vol.~3, Springer, 2006.
\bibitem[27]{sirignano2018dgm} {\sc J.~Sirignano and K.~Spiliopoulos}, {\em {DGM}: A deep learning algorithm for solving partial differential equations}, Journal of Computational Physics, 375 (2018), pp.~1339--1364.
\bibitem[28]{tsitsiklis2001regression} {\sc J.~Tsitsiklis and B.~Van~Roy}, {\em Regression methods for pricing complex {A}merican-style options}, IEEE Transactions on Neural Networks, 12 (2001), pp.~694--703.
\bibitem[29]{wang2018deep} {\sc H.~Wang, H.~Chen, A.~Sudjianto, R.~Liu, and Q.~Shen}, {\em Deep learning-based {BSDE} solver for {LIBOR} market model with application to {B}ermudan swaption pricing and hedging}, arXiv preprint arXiv:1807.06622, (2018).
\bibitem[30]{wang2009pricing} {\sc Y.~Wang and R.~Caflisch}, {\em Pricing and hedging {A}merican-style options: a simple simulation-based approach}, The Journal of Computational Finance, 13 (2009), pp.~95--125.
\bibitem[31]{yang2023optionpricing} {\sc J.~Yang and G.~Li}, {\em On sparse grid interpolation for {A}merican option pricing with multiple underlying assets}, arXiv preprint arXiv:2309.08287, (2023).
\bibitem[32]{zhang2017backward} {\sc J.~Zhang}, {\em Backward stochastic differential equations}, Springer, 2017.
\end{thebibliography}
\end{document}
